\section{Pressure tangents and the equilibrium magnetization radius}
\label{sec:pressure}

The variational principle bounds the magnetization of every state in
$\mathcal E_\beta$ by $m$.  We now prove that every ergodic member has
magnetization of length exactly $m$.  The obstruction to a smaller radius
will be analyticity of certain pressure tangents.  To obtain that
analyticity, we construct random transfer matrices whose traces are
nonzero before taking expectations; Lemma~\ref{lem:replica-cross} then
identifies the pressure tangents of their moments.

For an integer $r\ge2$, write $C_r=\{0,\ldots,r-1\}^3$, and let
$\mathcal C_r$ denote all its translates, on the lattice or on a sufficiently
large torus.  For any such cube $C$, put
$X_C=\sum_{x\in C}X_x$, and define $Y_C,Z_C$ in the same way.  Fix
$B\ge2$ and $W\ge2B$.  We will choose a constant $D_B>0$, independent of
$W$, and consider the nonpositive local observable
\begin{equation}
 a_{B,W}=-\frac{Y_{C_B}^2}{B^6}
         -D_B\frac{X_{C_W}^2+Z_{C_W}^2}{W^6}.
 \label{eq:negative-square-local}
\end{equation}
In an ergodic state with magnetization vector $(0,r,0)$, the transverse
term tends to zero as $W\to\infty$ at fixed $B$, while the $Y$ term
tends to $-r^2$ as $B\to\infty$, by
Lemma~\ref{lem:ergodic-variance}.
The periodic sum of $a_{B,W}$ is denoted by
$A_{B,W,L}=\sum_{x\in\Lambda_L}\tau_x(a_{B,W})$.  Define
\begin{equation}
 s_{B,W}(x)=\max_{\omega\in\mathcal E_\beta}
       \{\omega(a_{B,W})+x\omega(Y_0)\},\qquad x\in\R.
 \label{eq:square-tangent}
\end{equation}
This function is convex and $S$-Lipschitz.  It is even: a rotation by
$\pi$ around the $X$-axis preserves $a_{B,W}$ and sends $Y_0$ to $-Y_0$.

\begin{proposition}[Harmonic pressure tangent]
\label{prop:harmonic-tangent}
For every integer $B\ge2$ there is $D_B>0$ such that, for every integer
$W\ge2B$, the function in \eqref{eq:square-tangent} has a harmonic
extension $v_{B,W}$ from the positive real axis to
$\mathbb H=\{z\in\C:\Re z>0\}$ satisfying
\begin{equation}
 v_{B,W}(x)=s_{B,W}(x)\quad(x>0),\qquad
 v_{B,W}(z)\le s_{B,W}(\Re z)\quad(z\in\mathbb H).
 \label{eq:harmonic-tangent}
\end{equation}
\end{proposition}

We prove the proposition in the next three subsections.  We first
express the negative squares in
\eqref{eq:negative-square-local} as infinitesimal contributions of real
single-slot matrices, together with a small reduction of the
ferromagnetic coupling.  Shared Gaussian coefficients then turn moments
of the random traces into replicated pressures.  Finally a second-moment
argument selects nonvanishing traces whose logarithms converge to the
required harmonic function.

\subsection{A Laplacian construction for the negative squares}

Recall that there are $\ell=2S$ slots per site and that $P_*$ projects
onto the physical spin space.  For a real vector $v=(v_i)$ on finitely
many slots, write
\[
 D_v^a=\sum_i v_i s_i^a.
\]
Our generators will be $iD_v^2$ and $D_v^1,D_v^3$, with $\sum_i v_i=0$.
All three are real matrices in the $s^3$ basis, and their exponentials
with real coefficients are products of single-slot matrices in
$SL_2(\R)$.  Squaring $iD_v^2$ contributes a negative square in the
$Y$-direction.  The other two generators contribute positive squares;
the reduction of the static coupling will cancel their unwanted
nearest-neighbor coefficients.

All quadratic coefficients below use \emph{ordered} slot pairs:
\[
 (D_v^a)^2=\sum_{i,j}v_iv_j s_i^a s_j^a.
\]
For a cube $C$, let $L_C^{\rm full}$ be the graph Laplacian of the
complete graph on its $\ell|C|$ slots.  Let $L_C^{\rm bond}$ be the
Laplacian with an edge between every pair of slots whose sites are
nearest neighbors inside $C$.  Thus an off-diagonal Laplacian entry is
$-1$ on an edge.  These matrices are extended by zero when used on a
larger slot set.

The bond graph of $C_B$ is connected.  Consequently one can choose
$c_B>0$ so that, for every $C\in\mathcal C_B$,
\begin{equation}
 Q_C=c_B L_C^{\rm bond}-B^{-6}L_C^{\rm full}\ge0.
 \label{eq:small-cube-covariance}
\end{equation}
Indeed, both matrices annihilate the constant vector, and on its
orthogonal complement the bond Laplacian is positive definite while
$L_C^{\rm full}=\ell B^3I$.  Decompose $Q_C$ spectrally as
$\sum_v vv^{\mathsf T}$.  Every vector in this sum is real and has zero
sum; use the generators $iD_v^2$ for this list.

To cancel the remaining static reduction in the $X,Z$ directions, put
$n_{B,W}=(W-B+1)^3$ and, for $C\in\mathcal C_W$, consider
\begin{equation}
 R_C=D_BW^{-6}L_C^{\rm full}
       -\frac{c_B}{n_{B,W}}
          \sum_{\substack{C'\in\mathcal C_B\\C'\subseteq C}}
                      L_{C'}^{\rm bond}.
 \label{eq:large-cube-covariance}
\end{equation}
The constant $D_B$ can be chosen so that $R_C\ge0$ for every $W\ge2B$.
Here is a uniform bound.  Each slot belongs to at most $B^3$ small cubes,
and its degree in each bond graph is at most $6\ell$.  The sum of the
bond Laplacians therefore has norm at most $12\ell B^3$.  Since
$n_{B,W}\ge W^3/8$, the second term in
\eqref{eq:large-cube-covariance} has norm at most
$96c_B\ell B^3/W^3$.  On the zero-sum subspace, the first term is
$D_B\ell/W^3$ times the identity.  Thus, for example,
\begin{equation}
 D_B=96c_BB^3+1
 \label{eq:DB-choice}
\end{equation}
suffices.  Decompose each $R_C$ as $\sum_v vv^{\mathsf T}$ with real
zero-sum vectors, and use both generators $D_v^1$ and $D_v^3$.  Choose
the lists on one cube of each size and translate them.  The resulting
families are translation invariant.

Let $T_{B,L}$ be the physical compression of the static reduction
\begin{equation}
 -c_B\sum_{C\in\mathcal C_B}
     \sum_{\substack{i,j\text{ slots in }C\\
                     x_i,x_j\text{ nearest neighbors}}}
                   \sum_{a=1}^3s_i^a s_j^a,
 \label{eq:static-reduction}
\end{equation}
where the inner sum is ordered.  Each unordered site bond belongs to
$(B-1)B^2$ translates of $C_B$, so
\[
 H_{0,t}:=H_0+tT_{B,L}
   =\bigl(\beta-2tc_B(B-1)B^2\bigr)
                    \sum_{\{x,y\}\text{ bond}}\mathbf S_x\cdot\mathbf S_y.
\]
For fixed $B$ this remains ferromagnetic for all sufficiently small
$t>0$, uniformly in $L>2W$.

Denote all the generators just constructed by $G_\alpha$.  Their
quadratic contribution is compressed \emph{after squaring}:
\begin{equation}
 T_{B,L}+\sum_\alpha P_*G_\alpha^2P_*
       =A_{B,W,L}+C_{B,W}VI
 \label{eq:pulse-identity}
\end{equation}
on physical space, for a real scalar $C_{B,W}$ independent of $L>2W$.
To check the identity, first consider distinct slots $i,j$ in the cube
$C$ under consideration.  If $E_C(i,j)$
is the indicator of an edge of the bond graph, the coefficients from
a small cube in the $Y$-direction are
\[
 \underbrace{-c_BE_C(i,j)}_{\text{static reduction}}
 +\underbrace{\bigl(c_BE_C(i,j)-B^{-6}\bigr)}_{-(Q_C)_{ij}}
 =-B^{-6}.
\]
For either the $X$- or the $Z$-direction, the contribution from a large
cube is
\[
 (R_C)_{ij}=-D_BW^{-6}
       +\frac{c_B}{n_{B,W}}
          \sum_{\substack{C'\in\mathcal C_B\\C'\subseteq C}}E_{C'}(i,j).
\]
Every small cube lies in exactly $n_{B,W}$ large cubes.  After summing
over translates, the second term cancels the static reduction, leaving
the coefficient $-D_BW^{-6}$ on each ordered pair in each large cube.
These are precisely the off-diagonal coefficients of $A_{B,W,L}$.
The remaining terms have $i=j$ and are scalar because
$(s_i^a)^2=I/4$; translation invariance gives the scalar in
\eqref{eq:pulse-identity}.  In particular, the identity does not replace
$P_*G_\alpha^2P_*$ by $(P_*G_\alpha P_*)^2$, which would in general
give different single-replica terms.

\subsection{Nonvanishing random traces and replicated moments}

Fix $B,W$ and a sufficiently small $t>0$.  For an even $L>2W$ and an
integer $k\ge1$, take independent real centered Gaussian variables
$\lambda_{r,\alpha}$, $1\le r\le k$, of variance $2t/k$.  Fix an order
of the generators and define a slot-space layer
\[
 U_r(z)=\left(\prod_\alpha e^{\lambda_{r,\alpha}G_\alpha}\right)
                       e^{(tz/k)\sum_i s_i^2}.
\]
Only the complete layer is compressed.  With all products taken in
their fixed order, define the normalized random entire function
\begin{equation}
 \mathcal Z_{L,k,t}(z)=
 \frac{e^{-tC_{B,W}V}}{\Tr e^{H_0}}
 \Tr_{\rm phys}\prod_{r=1}^k
       \left[e^{H_{0,t}/k}\bigl(P_*U_r(z)P_*\bigr)\right].
 \label{eq:random-partition}
\end{equation}
Gaussian exponential moments ensure that all finite products of these
functions are integrable, locally uniformly in their arguments.

Every realization of \eqref{eq:random-partition} is nonzero on
$\mathbb H$.  Although the pulse supports overlap, their fixed ordered
product is still a tensor product of single-slot matrices: factors on
different slots commute and can be regrouped without changing the
order on any one slot.  All pulse matrices on a slot lie in
$SL_2(\R)$ and hence act as disk automorphisms in the $s^2$ basis, by
Lemma~\ref{lem:strip-maps}.  On every slot the remaining factor $e^{tzs^2/k}$ sends the
closed disk strictly inside the disk when $\Re z>0$.  Their composition
therefore has the strict mapping property of
Lemma~\ref{lem:trace-zero-free}.  To identify the permutation layers,
write $E$ for the full-slot exponential of the static interaction
with time step $1/k$.  It commutes with $P_*$, and cyclicity gives
\[
 \Tr_{\rm phys}\prod_{r=1}^k
       [e^{H_{0,t}/k}(P_*U_r(z)P_*)]
      =\Tr_{\rm slots}\prod_{r=1}^k[(EP_*)U_r(z)].
\]
Each $EP_*$ has nonnegative exchange rates and a symmetric seam, so
it is a permutation layer of the class allowed in that lemma.  The
identity uses $P_*^2=P_*$ and does not require $U_r$ to preserve
physical space.  This proves
samplewise nonvanishing, and in particular
\begin{equation}
 z\longmapsto\frac{1}{tV}\log|\mathcal Z_{L,k,t}(z)|
 \quad\text{is harmonic on }\mathbb H.
 \label{eq:sample-harmonic}
\end{equation}

Complex conjugation has a specific sign here.  The static matrices,
the projections and all pulses are real in the $s^3$ basis, whereas
$s^2$ is purely imaginary.  Therefore, for the \emph{same} Gaussian
sample,
\begin{equation}
 \overline{\mathcal Z_{L,k,t}(z)}
       =\mathcal Z_{L,k,t}(-\overline z).
 \label{eq:sample-conjugation}
\end{equation}
The argument on the right may lie outside $\mathbb H$; the function is
entire, so this causes no difficulty.  In particular its second moment
uses the pair $(z,-\overline z)$, not $(z,\overline z)$.

We next compute products of traces using replicas, always sharing the
Gaussian variables between replicas.  For fixed $j$ and arguments
$z_1,\ldots,z_j\in\C$, the finite-dimensional product formula gives
\begin{equation}
 \begin{split}
 \lim_{k\to\infty}\E\prod_{r=1}^j\mathcal Z_{L,k,t}(z_r)
  &=\mathcal R^{(j)}_{L,t}(z_1,\ldots,z_j)\\
  &:=\frac{\Tr\exp\!\left\{
        \sum_{r=1}^j\bigl(H_0^{[r]}+tA_{B,W,L}^{[r]}
                                      +tz_rM_L^{2,[r]}\bigr)
                         +tA_{{\rm cross},L}^{(j)}\right\}}
                       {(\Tr e^{H_0})^j}.
 \end{split}
 \label{eq:replicated-moment}
\end{equation}
The limit is locally uniform in the arguments.  To verify both the
limit and its cross terms, let $P^{(j)}$ be the product of the physical
projections in all replicas.  For one shared Gaussian variable,
\[
 \E\exp\!\left(\lambda\sum_{r=1}^jG_\alpha^{[r]}\right)
       =\exp\!\left(\frac tk
                          \left(\sum_{r=1}^jG_\alpha^{[r]}\right)^2\right).
\]
Independence allows these expectations to be taken in the prescribed
pulse order.  After the outer compression, the first-order term is
the following operator on physical replica space:
\begin{equation}
 P^{(j)}\left(\sum_{r=1}^jG_\alpha^{[r]}\right)^2P^{(j)}
   =\sum_{r=1}^j(P_*G_\alpha^2P_*)^{[r]}
       +2\sum_{r<s}(P_*G_\alpha P_*)^{[r]}
                         (P_*G_\alpha P_*)^{[s]}.
 \label{eq:compression-replicas}
\end{equation}
For an imaginary-direction generator $G_\alpha=iD_v^2$, the cross
product on the right is
\[
 -(P_*D_v^2P_*)^{[r]}(P_*D_v^2P_*)^{[s]}.
\]
For the other two directions
it has the positive sign.  Thus
\[
 A_{{\rm cross},L}^{(j)}
   =2\sum_\alpha\sum_{r<s}(P_*G_\alpha P_*)^{[r]}
                                      (P_*G_\alpha P_*)^{[s]}
\]
is self-adjoint, finite range and translation invariant.  Every one of
its local summands is a real multiple of a cross term in
Lemma~\ref{lem:replica-cross}, since all the vectors $v$ have zero sum.

Using \eqref{eq:pulse-identity} and distributing the scalar correction
in \eqref{eq:random-partition} among the $k$ steps, the expected
replicated step is $I+k^{-1}\mathcal H+O(k^{-2})$, where $\mathcal H$
is the exponent in \eqref{eq:replicated-moment}.  Here $B,W,j,t,L$ are
fixed and the error is locally uniform in $(z_1,\ldots,z_j)$.
Independence between time steps makes the expected product the $k$th
power of this deterministic step.  The matrix product limit
\cite{Trotter1959} proves \eqref{eq:replicated-moment}.

For real $x_1,\ldots,x_j$, the numerator defining
$\mathcal R^{(j)}_{L,t}$ is positive.  The thermodynamic limit followed
by the right pressure tangent gives
\begin{equation}
 \lim_{t\downarrow0}\lim_{L\to\infty}
       \frac{1}{tV}\log\mathcal R^{(j)}_{L,t}(x_1,\ldots,x_j)
           =\sum_{r=1}^j s_{B,W}(x_r).
 \label{eq:real-replica-tangent}
\end{equation}
Indeed, by \eqref{eq:pressure-tangent}, the left side is the maximum
of the perturbation density over equilibrium states of the unperturbed
replicated interaction.  Lemma~\ref{lem:replica-cross} makes the cross
density vanish in every such state.  Each marginal belongs to
$\mathcal E_\beta$, giving the upper bound by the sum on the right.
The product of maximizing single-replica states is an equilibrium
state and attains that sum.  This proves equality.

We also need an upper bound at complex arguments.  For self-adjoint
matrices $A,B$,
\begin{equation}
 |\Tr e^{A+iB}|\le\Tr e^A.
 \label{eq:complex-trace-bound}
\end{equation}
For completeness, apply the product formula to
$(e^{A/n}e^{iB/n})^n$ and then Schatten H\"older: its trace norm is at
most $\|e^{A/n}\|_n^n=\Tr e^A$, since the intervening factors are
unitary.  Taking the limit proves the inequality.  Applied to
\eqref{eq:replicated-moment}, it bounds its absolute value by
$\mathcal R^{(j)}_{L,t}(\Re z_1,\ldots,\Re z_j)$.

The limits in \eqref{eq:real-replica-tangent} are uniform on compact
sets of real arguments in their stated order.  Indeed the normalized
log traces are $S$-Lipschitz in each real argument, uniformly in $L,t$,
by the norm bound $\|M_L^2\|\le SV$; finite nets promote pointwise
convergence to uniform convergence.  The same bound holds for their
limits.  Combining this observation with
\eqref{eq:sample-conjugation}, \eqref{eq:replicated-moment} and
\eqref{eq:complex-trace-bound}, we obtain the following useful finite
parameter statement.  Given a compact $K\subset\mathbb H$, finitely
many positive real points $x_1,\ldots,x_q$, and $\epsilon>0$, one can
choose $t>0$ sufficiently small, then an even $L$ sufficiently large,
and finally $k$ sufficiently large, so that, with $T=tV$ and
$\mathcal Z=\mathcal Z_{L,k,t}$,
\begin{align}
 \E|\mathcal Z(z)|^2
      &\le \exp\{T(2s_{B,W}(\Re z)+\epsilon)\},\qquad z\in K,
          \label{eq:complex-second-moment}\\
 \left|T^{-1}\log\E U-2\sum_{i=1}^q s_{B,W}(x_i)\right|
      &\le\epsilon,
          \label{eq:product-first-moment}\\
 \left|T^{-1}\log\E U^2-4\sum_{i=1}^q s_{B,W}(x_i)\right|
      &\le\epsilon,
 \qquad U=\prod_{i=1}^q|\mathcal Z(x_i)|^2.
          \label{eq:product-second-moment}
\end{align}
The parameter $T$ can be required to exceed any prescribed number,
by enlarging $L$ after $t$ has been fixed.  For the last two estimates,
use the real argument lists $(x_1,-x_1,\ldots,x_q,-x_q)$ and twice that
list, together with the evenness of $s_{B,W}$.  These products are
nonnegative random variables at finite $k$.  Their limits at fixed
$t,L$ are strictly positive, so the locally uniform product limit also
gives the logarithmic estimates.  For the first estimate it gives the
required upper bound after the finite-volume parameters have been
fixed, even when its exponential right side is small.

This order of choices is essential: the matrix product limit is taken
at fixed volume; the pressure limit then precedes differentiation at
$t=0$.  The finite estimates are implemented in the reverse order,
choosing $t$, then $L$, then $k$.  No zero-free assertion about an
averaged partition function has been used.

\subsection{Selecting harmonic logarithms}

We isolate the analytic step to explain how moment information retains
the samplewise nonvanishing just proved.

\begin{lemma}[Selection from moment bounds]
\label{lem:harmonic-selection}
Let $s:(0,\infty)\to\R$ be locally Lipschitz, and enumerate a
countable dense subset $D=\{x_1,x_2,\ldots\}$ of $(0,\infty)$ with
$x_1=1$.  Suppose
$F_n$ are random holomorphic functions that are almost surely nonzero
on $\mathbb H$, and $T_n\to\infty$.  Assume, locally uniformly for
$z\in\mathbb H$, that
\begin{equation}
 \limsup_{n\to\infty}
   \left\{T_n^{-1}\log\E|F_n(z)|^2-2s(\Re z)\right\}\le0.
 \label{eq:selection-upper}
\end{equation}
For every $q\ge1$, suppose also that, with
$U_n=\prod_{i=1}^q|F_n(x_i)|^2$,
\begin{equation}
 T_n^{-1}\log\E U_n\longrightarrow2\sum_i s(x_i),\qquad
 T_n^{-1}\log\E U_n^2\longrightarrow4\sum_i s(x_i).
 \label{eq:selection-products}
\end{equation}
Then there is a harmonic function $v$ on $\mathbb H$ with
$v(x)=s(x)$ for every $x>0$ and $v(z)\le s(\Re z)$ everywhere.
\end{lemma}

\begin{proof}
Fix $q\ge1$, a compact $K\subset\mathbb H$ containing
$x_1,\ldots,x_q$, and $\delta>0$.  Put
$h_n=T_n^{-1}\log|F_n|$.
The Paley--Zygmund inequality \cite{PaleyZygmund1932} and
\eqref{eq:selection-products} give, for each $\epsilon>0$ and all
sufficiently large $n$,
\begin{equation}
 \PP\{U_n\ge\tfrac12\E U_n\}
   \ge\frac{(\E U_n)^2}{4\E U_n^2}
   \ge\tfrac14 e^{-3\epsilon T_n}.
 \label{eq:selection-paley}
\end{equation}
We compare this with the probability that $h_n$ exceeds
$s(\Re z)+\delta$ somewhere on $K$.

Cover $K$ by finitely many disks $D_j$ with concentric doubled disks
$\widetilde D_j$ compactly contained in $\mathbb H$, chosen so that
the oscillation of $s(\Re z)$ on each $\widetilde D_j$ is at most
$\delta/4$.  Write $s_j$ for its minimum there.  The submean
inequality for the nonnegative subharmonic function $|F_n|^2$ gives
\[
 \sup_{D_j}|F_n|^2
    \le C_j\int_{\widetilde D_j}|F_n(z)|^2\dd A(z).
\]
Using \eqref{eq:selection-upper}, the expectation of the right side
is at most $C'_j\exp\{T_n(2s_j+\delta/2+\epsilon)\}$ for large $n$.
If $h_n(z)>s(\Re z)+\delta$ at any point of $D_j$, then its left
side exceeds $\exp\{T_n(2s_j+2\delta)\}$.  Markov's inequality and a
sum over the disks therefore yield
\begin{equation}
 \PP\{h_n(z)>s(\Re z)+\delta\text{ for some }z\in K\}
       \le C_K e^{-(3\delta/2-\epsilon)T_n}.
 \label{eq:selection-bad-upper}
\end{equation}
Choose $\epsilon>0$ with $4\epsilon<3\delta/2$.  Since $T_n\to\infty$,
the bound in \eqref{eq:selection-bad-upper} is eventually smaller
than the lower bound in \eqref{eq:selection-paley}.  There is
therefore a realization for which both
\[
 h_n(z)\le s(\Re z)+\delta\quad(z\in K),\qquad
 \sum_{i=1}^q h_n(x_i)
     \ge\sum_{i=1}^q s(x_i)-\frac\epsilon2-
                                     \frac{\log2}{2T_n}
\]
hold.  Subtracting the upper bounds at the other $q-1$ points gives
\begin{equation}
 s(x_i)-(q-1)\delta-\frac\epsilon2-\frac{\log2}{2T_n}
       \le h_n(x_i)\le s(x_i)+\delta.
 \label{eq:selection-pointwise}
\end{equation}

Exhaust $\mathbb H$ by compact rectangles, and at stage $q$ enlarge
the rectangle to include the
first $q$ points.  Apply the preceding construction with
$\delta=q^{-2}$ and with $\epsilon$ small enough as above, choosing
an increasing sequence of indices.  The selected harmonic functions
are locally bounded above, and their values converge to $s$ at each
point of $D$.  Their values at $1$ in particular remain bounded
below.  On any connected relatively compact subdomain containing
$1$, subtracting these functions from a common upper bound and using
Harnack's inequality gives local lower bounds as well.  Interior
estimates and a diagonal subsequence yield a locally uniform harmonic
limit $v$.  The compact upper bounds give $v(z)\le s(\Re z)$; the
values on $D$, followed by continuity, give $v(x)=s(x)$ for every
$x>0$.
\end{proof}

\begin{proof}[Completion of the proof of Proposition~\ref{prop:harmonic-tangent}]
Fix $B,W$ and enumerate a countable dense set of positive real points,
starting with $1$.
Apply \eqref{eq:complex-second-moment}--\eqref{eq:product-second-moment}
successively on an exhaustion of $\mathbb H$, with errors tending to
zero and with $T_n=t_nL_n^3\ge n$.  At stage $n$ require the product
estimates for the first $q$ points for each $q\le n$.
This is a finite collection of requirements,
so the same parameter choices can enforce all of them.  The resulting
random functions $F_n=\mathcal Z_{L_n,k_n,t_n}$ satisfy
Lemma~\ref{lem:harmonic-selection}, by samplewise nonvanishing and the
moment estimates.  Its conclusion is precisely
\eqref{eq:harmonic-tangent}.
\end{proof}

\subsection{All ergodic equilibrium states have the same radius}

The negative square in \eqref{eq:negative-square-local} now separates
two hypothetical magnetization radii: a state with smaller radius
would dominate at small positive $x$, while the largest radius must
dominate at large $x$.  Harmonicity prevents this change.

\begin{proposition}[Equilibrium magnetization radius]
\label{prop:equilibrium-radius}
Every ergodic translation-invariant state $\omega\in\mathcal E_\beta$
satisfies
\begin{equation}
 \omega(X_0)^2+\omega(Y_0)^2+\omega(Z_0)^2=m^2.
 \label{eq:equilibrium-radius}
\end{equation}
\end{proposition}

\begin{proof}
The maximum of $\omega(Y_0)$ over the compact face
$\mathcal E_\beta$ is $m$.  The maximizing states form a nonempty
compact face, so it contains an extreme point $\omega_m$ of the
translation-invariant state space.  Thus $\omega_m$ is ergodic.  The
bound on the length of every equilibrium magnetization vector implies
that its vector is $(0,m,0)$.

Suppose an ergodic equilibrium state has magnetization length
$m'<m$.  Rotate it to a state $\omega_{m'}$ with vector $(0,m',0)$;
rotations preserve equilibrium and translation ergodicity.  By
Lemma~\ref{lem:ergodic-variance}, for $r=m,m'$,
\[
 \omega_r(Y_{C_B}^2/B^6)\longrightarrow r^2,
 \qquad
 \omega_r\bigl((X_{C_W}^2+Z_{C_W}^2)/W^6\bigr)
                                          \longrightarrow0.
\]
For each $B$, choose $W(B)\ge2B$ so large that the second expectation,
multiplied by $D_B$, is at most $1/B$ in both states.  Then, writing
$s_B=s_{B,W(B)}$,
\begin{equation}
 \liminf_{B\to\infty}s_B(x)
        \ge\max\{xm-m^2,\;xm'-(m')^2\},\qquad x>0.
 \label{eq:two-radius-lower}
\end{equation}

There is also an upper bound that is uniform in $B,W$.  If
$y=\omega(Y_0)$, positivity of variance and translation invariance
give
\[
 \omega(Y_{C_B}^2/B^6)\ge y^2,\qquad -m\le y\le m.
\]
Dropping the other nonpositive square terms shows that
\begin{equation}
 s_{B,W}(x)\le\max_{-m\le y\le m}(xy-y^2)
                 =xm-m^2,\qquad x>2S,
 \label{eq:large-field-tangent}
\end{equation}
because $m\le S$ and $y\mapsto xy-y^2$ is increasing on $[-m,m]$
for such $x$.  For every $x>0$ we also have $s_{B,W}(x)\le Sx$.

Let $v_B$ be the harmonic function supplied by
Proposition~\ref{prop:harmonic-tangent}.  It satisfies
$v_B(z)\le S\Re z$ on $\mathbb H$.  At any fixed $x_0>2S$,
\eqref{eq:two-radius-lower} and \eqref{eq:large-field-tangent} give
$v_B(x_0)\to x_0m-m^2$.  Harnack's inequality applied to the
nonnegative harmonic functions $S\Re z-v_B(z)$, followed by interior
compactness, supplies a subsequential harmonic limit $v$ on
$\mathbb H$.  The same two inequalities imply
$v(x)=xm-m^2$ for every $x>2S$.  The restriction of a harmonic
function to the real axis is real analytic, so this identity holds
for every $x>0$.

On the other hand, \eqref{eq:two-radius-lower} passes to this
subsequence and gives $v(x)\ge xm'-(m')^2$.  For
$0<x<m+m'$, these two assertions contradict
\[
 \bigl(xm'-(m')^2\bigr)-\bigl(xm-m^2\bigr)
           =(m-m')(m+m'-x)>0.
\]
Since $m>0$, this interval is nonempty even when $m'=0$.
No smaller radius is possible, proving \eqref{eq:equilibrium-radius}.
\end{proof}
